\section*{अध्याय \ref{ch:g9:periodic-table-first-look} --- आवर्त सारणी: पहली झलक}

\begin{solution}{exo:g9:periodic-table-first-look:1}
बढ़ते \omterm{def:g9:inside-the-atom:atomic-number}{परमाणु क्रमांक} $Z$ के क्रम में।
\end{solution}

\begin{solution}{exo:g9:periodic-table-first-look:2}
कार्बन: \omterm{def:g9:periodic-table-first-look:period}{आवर्त} 2, \omterm{def:g9:periodic-table-first-look:period}{समूह} 14। मैग्नीशियम: \omterm{def:g9:periodic-table-first-look:period}{आवर्त} 3, \omterm{def:g9:periodic-table-first-look:period}{समूह} 2। आर्गन:
\omterm{def:g9:periodic-table-first-look:period}{आवर्त} 3, \omterm{def:g9:periodic-table-first-look:period}{समूह} 18।
\end{solution}

\begin{solution}{exo:g9:periodic-table-first-look:3}
\omterm{def:g9:periodic-table-first-look:metal}{धातुएँ}: लोहा, ताँबा, ऐलुमिनियम। \omterm{def:g9:periodic-table-first-look:metal}{अधातुएँ}: सल्फ़र, ऑक्सीजन, क्लोरीन।
\end{solution}

\begin{solution}{exo:g9:periodic-table-first-look:4}
सोडियम के नीचे पोटैशियम (\ce{K}); क्लोरीन के नीचे ब्रोमीन (\ce{Br})।
\end{solution}

\begin{solution}{exo:g9:periodic-table-first-look:5}
कैल्सियम (\ce{Ca}, $Z = 20$) और नाइट्रोजन (\ce{N}, $Z = 7$)।
\end{solution}

\begin{solution}{exo:g9:periodic-table-first-look:6}
वह सोडियम वाले समूह, \omterm{def:g9:periodic-table-first-look:period}{समूह} 1, में है, और एक समूह के तत्व एक जैसा व्यवहार
करते हैं।
\end{solution}

\begin{solution}{exo:g9:periodic-table-first-look:7}
\omterm{def:g9:periodic-table-first-look:metal}{धातु}: वह चमकीला है, विद्युत का चालन करता है और धनात्मक \omterm{def:g9:ions:ion}{आयन} बनाता है।
\end{solution}

\begin{solution}{exo:g9:periodic-table-first-look:8}
मिलते-जुलते गुणों वाले तत्वों को एक ही स्तंभ में रखने के लिए, भले ही इसके लिए
कोई खाना ख़ाली छोड़ना पड़े। ख़ाली जगहें भविष्यवाणियाँ थीं: लुप्त तत्व बाद में उन्हीं
गुणों के साथ मिले जिनका उन्होंने वर्णन किया था।
\end{solution}

\begin{solution}{exo:g9:periodic-table-first-look:9}
कार्बन: $\ce{C} = 12$; ऑक्सीजन: $\ce{O} = 16$। ख़ाली जगहें ``? $= 68$''
(Al $= 27.4$ के बगल में) और ``? $= 70$'' (Si $= 28$ के बगल में)।
\end{solution}

\begin{solution}{exo:g9:periodic-table-first-look:10}
\omterm{def:g9:periodic-table-first-look:period}{आवर्त} 2: 8 तत्व (लिथियम से नियॉन तक)। \omterm{def:g9:periodic-table-first-look:period}{आवर्त} 4: 18 तत्व (पोटैशियम से
क्रिप्टॉन तक)।
\end{solution}

\begin{solution}{exo:g9:periodic-table-first-look:11}
\omterm{def:g9:periodic-table-first-look:period}{समूह} 18। जहाँ कुछ भी अभिक्रिया नहीं करना चाहिए: बल्बों के भीतर और वेल्डिंग गैस
में आर्गन, गुब्बारों में हीलियम (वह \omterm{def:g4:what-a-fire-needs:burning}{जलती} नहीं)।
\end{solution}

\begin{solution}{exo:g9:periodic-table-first-look:12}
$(27.0 + 114.8)/2 = 70.9$, जो 69.7 के क़रीब है।
\end{solution}

\begin{solution}{pb:g9:periodic-table-first-look:1}
\textbf{1.} $Z = 32$; \omterm{def:g9:periodic-table-first-look:period}{आवर्त} 4, \omterm{def:g9:periodic-table-first-look:period}{समूह} 14।

\textbf{2.} ऊपर: सिलिकॉन। नीचे: टिन। बाएँ: गैलियम। दाएँ: आर्सेनिक।

\textbf{3.} उपधातु।

\textbf{4.} वह सिलिकॉन वाले समूह में है, और एक समूह के तत्व एक जैसा
व्यवहार करते हैं।

\textbf{5.} $(28.1 + 118.7)/2 = 73.4$।

\textbf{6.} $(69.7 + 74.9)/2 = 72.3$।

\textbf{7.} बाएँ--दाएँ का औसत (72.3): \omterm{def:g9:periodic-table-first-look:period}{आवर्त} में \omterm{def:g7:atoms-and-molecules:atom}{परमाणु} भार तत्व-दर-तत्व
लगातार बढ़ता है, जबकि समूह में नीचे जाने पर वह बड़ी छलाँगों में
बढ़ता है।

\textbf{8.} $72.6 - 70 = 2.6$।

\textbf{9.} $(28.1 + 118.7 + 69.7 + 74.9)/4 = 291.4/4 = 72.85$, यानी
72.9।

\textbf{10.} औसत, 72.9, मेंडेलीव के 72 और विंकलर के 72.3 (और आज के
72.6), दोनों के क़रीब है।

\textbf{11.} किसी के देखने से पहले ही जिस तत्व की भविष्यवाणी की गई थी, वह बताए गए
गुणों के साथ मिला: सारणी केवल एक सूची नहीं थी, वह भविष्यवाणियाँ
कर सकती थी।

\textbf{12.} लगभग 72.9।
\end{solution}
